% !TEX root = ../main.tex
\chapter{Form Invariance and Repeatability Both Fail}
\label{ch:canonical-form}
\label{sec:canonical-form}

One qualification belongs in the first paragraph, because the heading is stronger than
what follows and the difference is the point. Training on paraphrase-rich text does
apply pressure toward treating equivalent formulations alike, and embeddings do place
many paraphrases near one another. What is absent is not similarity; it is a
\emph{guarantee}---an enforced, inspectable, versionable canonical form of the kind
every other component in the stack supplies. A learned tendency and an enforced
invariant differ in exactly the way that matters for validation, and the rest of this
section is about that difference.

One property does most of the damage, and it is worth naming before it is formalized.
Call it \emph{form invariance}: materially equivalent requests should produce equivalent
decision-relevant outcomes. It is the property every other component in an enterprise
stack supplies without being asked, it is the property this architecture does not supply,
and its absence is what breaks validation. It is stated as
invariant I1 in Chapter~\ref{ch:invariants}, where the full set is tabulated and
where the numbering used throughout is defined; this section explains why it fails, what follows for a
validation program, and why the obvious remedy does not work. The failure has a
one-sentence statement: the system is keyed on the surface form of its input, not on the
content.

\section{The map and what it does not respect}

Section~\ref{sec:what-computed} established the signature: $S$ maps
$\mathcal{V}^{*}$ to distributions over $\mathcal{V}^{*}$. The domain is token
sequences.

Now consider the equivalence relation that validation cares about. Write
$x \sim_c x'$ when $x$ and $x'$ encode \emph{materially equivalent} requests under
the same resolved context $c$---same question, same governing frame, same material
facts, same intended use. This relation partitions the input space. What a test
suite needs is that decision-relevant behavior be constant on those classes:
\begin{equation}
  x \sim_c x' \;\Longrightarrow\; \text{EquivalentDecisionRelevantOutcome}(x, x', c).
  \label{eq:form-invariance}
\end{equation}

Two things about \eqref{eq:form-invariance} are deliberate. It is stated over
\emph{material} equivalence under a resolved context, not over semantic
equivalence in the abstract, because the equivalence that matters is domain- and
context-relative. And it requires equivalent decision-relevant \emph{outcomes},
not identical text. Nobody needs the system to produce the same sentence twice.
What an institution needs is that materially equivalent requests not yield
incompatible assertions, incompatible evidence bases, or incompatible
decision-relevant conclusions without an identifiable reason.

Nothing in the architecture delivers \eqref{eq:form-invariance}. There is no
normalization pass, no canonical form, no quotient map. By
Section~\ref{sec:what-computed} the tokenizer is lexical, the embedding is of
tokens, and no later stage brings two encodings of one claim to a common form. Two
members of a $\sim_c$ class are two different points in $\mathbb{R}^{n \times d}$
with different distributional histories, and \eqref{eq:objective} contains no term
penalizing divergence between them. The map is free to send them anywhere, and
sometimes does.

\section{Every other component in the stack does better}

The failure is easy to overlook because it is unlike anything else in an enterprise
architecture.

\begin{itemize}
  \item A database normalizes before it compares. Results do not depend on whether
        a predicate was written \texttt{NOT (a OR b)} or \texttt{(NOT a) AND (NOT
        b)}; the planner canonicalizes, and equivalence is decidable in the
        algebra. The comparison should not be overdrawn: equivalence for the full
        query language is not decidable either, and floating-point aggregation can
        differ across parallel plans. What a database supplies is a canonical form
        for the fragment that matters and a documented statement of where determinism
        is and is not guaranteed. Neither is available here.
  \item A compiler canonicalizes before it optimizes. \texttt{x*2} and
        \texttt{x+x} converge.
  \item A cache keyed correctly hashes normalized content, not raw bytes. Keying a
        cache on the unnormalized string is a familiar bug, and its symptom is
        familiar too: identical requests, different answers, no error anywhere.
\end{itemize}

The last case is the closest, and it is barely an analogy. The system behaves like a
cache keyed on the raw string: materially identical requests are distinct keys, and
nothing in the pipeline would notice.

The deeper version of the point is architectural, and it is the one that should
concern anyone who has built a system of record. A canonical record requires stable
identity, a normalized or governed representation, declared dependencies, lineage,
and reproducible derivation. A model's latent representation has none of these
properties: it has no stable identity for an entity across paraphrases, no
inspectable normalized form, no declared dependency on a source, and no reproducible
derivation path from evidence to output. It is therefore not a candidate system of
record, whatever its accuracy. Few architects would assert that it is; the error is
subtler and commoner, and consists in treating model \emph{output} as authoritative
without requiring the lineage that any other authoritative source in the enterprise
would have to supply.

Insofar as materially equivalent formulations \emph{are} treated alike---and often
they are---that is a contingent achievement of training coverage. The training
distribution contained many paraphrase pairs, so the map is approximately constant
on well-covered classes. This degrades exactly where prediction says it should: as
formulations become rarer, longer, more domain-specific, more like the queries a
specialist actually asks. The behavior is best where the guarantee is needed
least.

\section{What the measurements show}

The prediction is testable and has been tested.

The cleanest case is format sensitivity. Holding semantic content fixed
and varying only formatting choices---separators, casing, spacing, the shape of an
enumeration---produces large spreads in task performance \citep{sclar2024}. This is
the result to show a sceptical colleague: the content of the request is unchanged by
construction, the variations are choices a prompt author makes arbitrarily and does
not record, and behavior moves anyway. Aggregate performance cannot shift unless
individual answers changed, so a spread across formats is a lower bound on the rate
at which materially equivalent requests receive different treatment.

Framing sensitivity gives the direction, and is the second clean case. When a prompt signals what the user
would prefer to hear, assertion shifts toward it \citep{sharma2024}. The content
of the question is unchanged; the assertion is not.

Consistency under paraphrase, measured directly, falls well below what a
content-directed process would produce \citep{elazar2021}, and models will
contradict themselves within a single generation \citep{mundler2024}. The second
of these bears on repeatability as well as on form invariance: the instability is not only across sessions
and phrasings but inside one continuous output, where any notion of a maintained
position would have to apply if it applied anywhere.

\section{Why ``train a normalizer'' is not the answer}
\label{sec:normalizer}

The natural engineering response is to canonicalize the input: put a normalizer in
front of the model and the problem goes away. It does not, for five reasons that
compound.

\paragraph{Material equivalence is context-sensitive.} Whether two requests are
equivalent depends on $c$. ``Is this reportable?'' asked of two transactions
identical but for jurisdiction are not materially equivalent requests, and no
context-free normalizer can know that.

\paragraph{It depends on domain ontology.} Deciding that two phrasings pick out the
same product class, the same patient population, or the same instrument type
requires the domain's own identity conditions. That is ontology work, not string
work.

\paragraph{Paraphrases carry material distinctions.} In regulated language, surface
differences are frequently the point. ``Indicated for'' and ``may be used in'' are
not paraphrases. ``Shall'' and ``should'' are not paraphrases. A normalizer
aggressive enough to collapse genuine paraphrases will collapse these too, and the
failure will be silent and consequential.

\paragraph{General semantic equivalence is not reliably decidable.} Unlike
relational algebra, natural language offers no decision procedure for equivalence,
and any learned approximation inherits the properties of the thing it approximates,
including sensitivity to form.

\paragraph{A latent representation is not an inspectable canonical record.} Even a
normalizer that worked would produce an internal representation that no reviewer
can read, version, or diff---which fails the requirement it was introduced to
satisfy.

The conclusion is not that normalization is useless. Structured intake, controlled
vocabularies, and typed request objects all help, substantially, and belong in the
architecture. The conclusion is that the remedy is not \emph{semantic}
normalization of free text. It is explicit context binding, structured domain
representations where they can be built, governed source corpora, deterministic or
inspectable policy controls where the decision permits them, and a preserved
evidence-and-policy path for every consequential assertion. That is what governed
retrieval is for.

\section{The validation consequence}
\label{sec:validation-consequence}

Now the argument this book exists to make. State the validation inference
explicitly and it fails visibly.

Let $T = \{(x_1, a_1), \ldots, (x_n, a_n)\}$ be a test suite: inputs paired with
acceptable outputs. Suppose $S$ passes---for each $i$, $S(x_i)$ is acceptable. The
inference validation needs is:

\begin{quote}
For every $x$ that a user might actually send, if $x \sim x_i$ for some tested
$x_i$, then $S(x)$ is acceptable.
\end{quote}

That inference is valid only if \eqref{eq:signature} is constant on $\sim$-classes,
which is form
invariance. Form invariance fails. So the inference is invalid, and the test suite
does not license the conclusion drawn from it.

Three aggravating features, each of which an engineer will recognize as worse than
the bare failure.

\paragraph{The classes are unbounded.} You cannot enumerate the paraphrases of a
question, so you cannot fix the problem by testing more members of each class.
Coverage is not a strategy against an infinite equivalence class; the best you can
do is sample it, which returns you to sampling.

\paragraph{The failures are silent.} A form-invariance violation raises nothing. There is no
exception, no warning, no anomaly in the logs. The system returns a fluent,
well-formed, differently-wrong answer with the same confidence as the right one.
Contrast a schema violation or a type error, which announce themselves.

\section{The same failure, on a single input held fixed}
\label{sec:repeatability-fails}

Everything so far concerns \emph{different} inputs that mean the same thing. The
second failure is worse, because it does not need two inputs.

Repeatability --- invariant I7 --- is the property that the same request, in the same
session, yields the same assertion. It does not hold, and by
Section~\ref{sec:sampling} it cannot be assumed to: the output is a draw, and the
sequence conditioned on grows as the session does, so the request that was tested and
the request as it later arrives are not the same point in input space even when the
characters are identical.

The consequence is worth stating separately from the paraphrase argument, because it
is a strictly stronger claim. Form invariance failing means a passing test does not
\emph{extend} to materially identical requests --- the guarantee has a boundary, and
the boundary is narrower than anyone assumed. Repeatability failing means the
guarantee does not securely cover \emph{the string that was tested}. There is no
boundary to find, because the test result was a statement about one draw and the
next draw is a different event.

For a validation program these compose badly. Enumerate more paraphrases and you
address form invariance without touching repeatability. Re-run the same case a
hundred times and you estimate a per-input distribution without learning anything
about the neighboring inputs real users will send. Neither remedy reaches the other
failure, and no amount of either produces the thing a deployment decision wants,
which is a claim about behavior rather than about a sample of it.

\section{Why this makes accuracy the wrong thing to certify}

Put the two halves together.

Accuracy is estimated by sampling: draw inputs, compare outputs to expected
answers, report the proportion correct. The estimate is a statement about the
distribution the inputs were drawn from. Using it to certify deployment requires
that the estimate transfer to the inputs that will actually arrive---and the
failure of form invariance is precisely the claim that it does not transfer across
semantic
equivalence, which is the dimension along which real inputs differ from test
inputs.

The conclusion should be stated carefully, because the strong form is indefensible and
invites dismissal. Reference-based accuracy is not worthless and it does not establish
nothing. It provides evidence \emph{conditional on} a specified evaluation
distribution, task operationalization, prompt representation, model version, and serving
configuration. Every empirical validation regime generalizes from finite samples under
such conditions, and a classifier is in the same position with respect to distribution
shift unless its input domain and robustness properties are bounded.

The claim is narrower and it is about which conditions can be specified. In free-form
consequential interaction, the prompt representation is chosen by the user and is not
recorded; the equivalence class of materially identical requests is unbounded and cannot
be enumerated; and the map is not constant on that class. So the conditions under which
the accuracy estimate would transfer are not merely unmet but unstateable, and the
generalization that broad consequential reliance requires is unavailable --- not because
the measurement is wrong, but because what it is conditional on cannot be pinned down.
The number is a correct measurement of how the system performed on those strings, and
that is a smaller thing than it is usually taken for.

Notice, finally, what this does to the cell-(ii) case of
Section~\ref{sec:cell-ii}. A lucky assertion is one that would have been produced
whatever the fact had been. Verifying it tells you about the draw, not the
process. The failure of form invariance generalizes that observation from lucky assertions to
\emph{all} assertions: any verified output tells you about the draw. Cell (ii) was
the special case where the disconnection is total; form-invariance failure is the
same disconnection showing up as instability across the input space, and
non-repeatability is it showing up on a single input held fixed.

That is the link between the two halves of this book. The taxonomy of
Section~\ref{sec:taxonomy} says that a correct answer can be produced for no
reason. This section says that even a correct answer produced for a good reason
carries no guarantee of recurrence, because nothing holds the system to the
content rather than the string. Both point at the same missing property, and the
next section names it.

\section{Why the two claims are one}
\label{sec:unity}

Section~\ref{sec:taxonomy} redefined hallucination as ungrounded assertion. This
section argued that there is no canonical form, so a passing test establishes
nothing about a materially equivalent input. The two are the same architectural
fact seen from two sides, and it is worth saying how, because each makes the other
consequential.

Why do these systems hallucinate in the redefined sense? Because assertion
probability is computed over formulations. A string can be made probable by surface
statistics with no route through evidential support for its content, so assertion
and evidence come apart---which is exactly $\neg\,\Grd_S(p \mid c)$ while
$\Assert_S(p)$ holds. Ungrounded assertion is what the absence of a content-level
representation looks like from the output side.

And why does the absence of a canonical form matter, rather than being an
interesting curiosity about paraphrase robustness? Because of what
Section~\ref{sec:cell-ii} showed about a lucky assertion: verifying it tells you
about the draw, not the process. Form-invariance failure generalizes that from
lucky assertions to \emph{all} assertions. Any verified output tells you about the
draw, because nothing holds the system to the content rather than the string. The
lucky cell was the special case where the disconnection is total; form-invariance
failure is the same disconnection showing up as instability across the input space.

So the redefinition identifies the defect and the canonical-form argument explains
why you cannot test your way out of it. Together they force one conclusion: the property to measure is not whether assertions turn
out right, but whether they were grounded when made.
